\documentclass[11pt]{article}
\usepackage[a4paper,margin=27mm]{geometry}
\usepackage{amsmath,amssymb,amsthm,hyperref}
\newtheorem{theorem}{Theorem}
\newtheorem{lemma}[theorem]{Lemma}
\newcommand{\gen}[1]{\langle #1\rangle}
\newcommand{\CG}{\mathrm C}
\title{Extending a marked group by a centralizing involution}
\author{human\_compression project}
\date{11 October 2026}
\begin{document}
\maketitle

\section{The extension problem}
A string tuple of length $n$ in a group $G$ is an ordered tuple
$(r_0,\ldots,r_{n-1})$ of distinct nonidentity involutions such that
$r_i r_j=r_j r_i$ whenever $|i-j|>1$. It is a generating string marking
when its entries generate $G$. No intersection property is required.
The following argument isolates the group-theoretic step used in the
half-turn constructions of Leemans--Toledo \cite[Theorem~4.4 and
Proposition~4.7]{LT}, and in the corrected symmetric-end construction
of the reviewed application \cite[Sections~5--7]{Application}.

The inputs from projective geometry remain separate. In the application,
an even cyclic torus supplies an involution which centralizes an odd
rotation dihedral subgroup. An exact maximal-subgroup theorem supplies
the proper alternating or symmetric end. Neither existence statement is
proved by the certificate below. Its conclusion concerns the actual four
ordered generators, including their distinctness and generation.

Let $(a,b,c)$ be a string tuple in $G$, and put
\[
 H=\gen{a,b,c},\qquad D=\gen{a,b}.
\]
We use $\CG_H(a,b)=\{x\in H:xa=ax,\ xb=bx\}$; equivalently this is
$H\cap\CG_G(D)$. A maximal proper subgroup means an actual maximal
member of the subgroup lattice below $G$, not merely a subgroup maximal
within a selected list.

\begin{lemma}\label{lem:tuple}
Suppose that $z$ is a nonidentity involution commuting with $a$ and $b$,
and that $z$ differs from $a,b,c$. Then $(a,b,c,z)$ is a string tuple.
Its first three entries generate exactly $H$.
\end{lemma}
\begin{proof}
The old entries satisfy their original involution and distinctness
conditions, and the hypotheses supply these conditions for the new
entry. The nonconsecutive pairs of positions in a tuple of length four
are $(0,2),(0,3),(1,3)$. The corresponding commutations are $ac=ca$,
$az=za$ and $bz=zb$, all of which are hypotheses. Restricting the new
tuple to positions $0,1,2$ gives the original tuple entry by entry;
therefore its generated subgroup is $H$.
\end{proof}

\begin{lemma}\label{lem:outside}
If $\CG_H(a,b)=1$, every nonidentity element $z$ commuting with $a$ and
$b$ lies outside $H$.
\end{lemma}
\begin{proof}
If $z\in H$, the two commutations put it in $\CG_H(a,b)$. The
triviality assumption then gives $z=1$, a contradiction. This argument
does not require that $z$ have finite order.
\end{proof}

\begin{lemma}\label{lem:generation}
If $H$ is a maximal proper subgroup of $G$ and $z\notin H$, then
$\gen{H,z}=G$.
\end{lemma}
\begin{proof}
The subgroup $\gen{H,z}$ contains $H$. Its containment is strict because
it contains $z$, whereas $H$ does not. Maximality leaves $G$ as the only
possibility.
\end{proof}

\begin{theorem}[Full-end extension]\label{thm:full}
Assume $H=G$, and let $z$ be a nonidentity involution commuting with
$a,b$. Either of the following additional conditions suffices to make
$(a,b,c,z)$ a generating string marking of $G$:
\begin{enumerate}
\item $ab\ne ba$ and $bc\ne cb$;
\item $z\notin D$ and $z\ne c$.
\end{enumerate}
The subgroup on its first three generators remains $G$.
\end{theorem}
\begin{proof}
Under the first condition, $z=a$ would imply $ab=ba$ by the commutation
of $z$ with $b$. The equality $z=b$ would similarly imply $ab=ba$ by
commutation with $a$. Finally, $z=c$ would imply $bc=cb$. These three
contradictions prove the required distinctness. Under the second
condition, $a,b\in D$ immediately exclude the first two coincidences,
and the last is explicitly excluded. Lemma~\ref{lem:tuple} now supplies
all string axioms. The new generated subgroup contains $H=G$, so it is
$G$. The prefix is unchanged.
\end{proof}

\begin{theorem}[Maximal-end extension]\label{thm:maximal}
Assume that $H$ is a maximal proper subgroup of $G$. Let $z$ be a
nonidentity involution commuting with $a,b$. If $z\notin H$, then
$(a,b,c,z)$ is a generating string marking of $G$, with first-three
subgroup exactly $H$. In particular, the same conclusion follows from
$\CG_H(a,b)=1$ without assuming $z\notin H$ separately.
\end{theorem}
\begin{proof}
All three old generators belong to $H$. Thus $z\notin H$ makes $z$
different from each of them. Lemma~\ref{lem:tuple} proves the string
axioms, and Lemma~\ref{lem:generation} proves generation. If the internal
common centralizer is trivial, Lemma~\ref{lem:outside} first proves the
required nonmembership. Prefix equality again identifies the actual
old subgroup; this conclusion is stronger than an abstract isomorphism
of unnamed subgroups.
\end{proof}

\section{Relation to the projective-line application}
For the full-end case in \cite[Section~5]{Application}, let $q$ be odd,
choose $\alpha\in\{1,-1\}$ with $q\equiv\alpha\pmod4$, and put
$N=(q-\alpha)/2$. The application retains the usual split and nonsplit
tori of $\operatorname{PSL}_2(q)$ and their normalizers. If the first
adjacent order is an odd divisor $k>1$ of $N$, the relevant torus has
order $N$, which is even. Its unique involution $z$ centralizes the
normalizer and hence $D$. Since $D$ has odd rotation order and trivial
centre, $z\notin D$. The next adjacent pair is noncommuting, so $z\ne c$.
Theorem~\ref{thm:full}(2) is the final extension step. Alternatively,
Theorem~\ref{thm:full}(1) uses the already proved noncommutation of both
adjacent pairs and avoids repeating the distinctness argument.

For the proper-end cases, the accepted application proves
$\CG_H(a,b)=1$ for its chosen $D_{10}$ in $A_5$ and $S_3$ in $S_4$.
It retains Giudici's exact maximal-subgroup theorem in the applicable
prime ranges \cite[Theorem~2.2(6),(8)]{Giudici}. The torus still supplies
$z$, and Theorem~\ref{thm:maximal} then proves generation and the actual
facet subgroup. The following steps remain in the written application:
the opposite end is full by the classification; the marked-group action
produces a regular maniplex; subgroup indices give the face counts;
and simplicity gives nonorientability. None is silently folded into
the new Lean conclusion.

\section{Formal interface and limitations}
The Lean file defines local string tuples and generating string markings
using an actual function $\mathrm{Fin}(n)\to G$. Its conclusions contain
nonidentity, involution equations, function injectivity, every distant
commutation and subgroup-closure equality with the whole group.
The concrete tuple is $[a,b,c,z]$. The theorem
\texttt{prefix\_preserved} identifies its restriction with $[a,b,c]$.
The internal-centralizer hypothesis is the subgroup equality
\[
 H\cap \CG_G(\{a,b\})=\{1\},
\]
represented by an infimum and the bottom subgroup. Maximality is
\texttt{IsCoatom H} in the full subgroup lattice.

The Lean certificate proves the lemmas and extension theorems above
for arbitrary groups. It assumes neither the four-generator conclusion
nor the desired final generation statement. It does assume the stated
old triple, and the projective applications must supply its hypotheses.
There is no formalization here of projective groups, torus existence,
Dickson or Giudici classification, the finite local census, the
intersection property, maniplexes, or the final infinite-family
statements. The compiled declarations use only Lean's standard logical
axioms. Library dependencies and their exact revisions are recorded in
the replay package; the imported compiled library remains trusted.

An elementary finite control explains why the outside condition is
substantive even if distinctness is independently known. In the additive
group $G=\mathbb F_2^4$, write products additively and take
$a=e_1$, $b=e_2$, $c=e_3$, $z=e_1+e_2+e_3$. The subgroup
$H=\gen{a,b,c}$ is a maximal hyperplane. All four entries are distinct
nonzero involutions and all commute, but they still generate $H$,
not $G$. Here $z\in H$ and $\CG_H(a,b)=H$ is nontrivial.
The companion finite checker exhausts the sixteen elements and verifies
this countermodel, including maximality by testing every outside
adjoining element. It is an independent finite check of the weakened
statement's failure, not the proof of the general theorems.

\begin{thebibliography}{9}
\bibitem{LT} D.~Leemans and M.~Toledo,
\emph{Maniplexes with automorphism group $\operatorname{PSL}(2,q)$},
author manuscript dated 15 May 2023, eleven pages,
Theorem~4.4 and Proposition~4.7.
\bibitem{Application} human\_compression project,
\emph{Centralizer extensions of projective-line maniplexes},
reviewed phase258 application, Sections~5--7; independent review phase263.
\bibitem{Giudici} M.~Giudici,
\emph{Maximal subgroups of almost simple groups with socle
$\operatorname{PSL}(2,q)$}, arXiv:math/0703685v1, 2007,
Theorem~2.2.
\end{thebibliography}
\end{document}
